\section{Method: building an environment}\label{method-building-an-environment}

The method has two parts. The first turns what the creator wants into a locked design: a concept for the environment and for each of its scenes, an inventory of every object, and a plan in metres. The second reads the layers each scene needs off that design and makes each in turn, behind its gates, until the scene passes its completion gate.

\subsection{From concept to plan}\label{from-concept-to-plan}

Step 1: the environment concept. The creator first locks the EC in a short interview with Gemini 3.1 Pro \citep{gemini31pro}, in which he fixes its style and its physics, adds his reference pictures and names the scenes it holds. An environment plan then places each top scene's site on that ground and works out the sky each site sees.

Step 2: a concept for each scene. A concept is then locked for every scene, in the order of parent scene before child scene. Each draw is one Nano Banana Pro \citep{nanobananapro} call, made on a concept review page where the creator can mark areas and ask for a new draw. A child scene is drawn only once its parent has been locked and planned, and it is drawn on a crop of the parent's concept, so that it fits the parent it stands in.

Step 3: inventory and plan. Gemini 3.1 Pro lists every object of the locked concept as a row of the scene inventory, giving its kind, count, anchor and size in metres, and SAM 3 \citep{carion2025sam3} draws a box around each row; nothing that is not a row is made or placed. Nano Banana Pro then draws the scene from straight above, and Gemini 3.1 Pro writes the plan in metres from that drawing, the concept and the inventory. The plan gives the outline, walls, openings, levels, stairs, roofs and terrain, each item with its facing, people as the spots on which they stand, and the rooms of the child scenes. A set of free checks follows: every row must be placed as often as its count says, lie inside the outline and not overlap another floor item, every door and standing spot must be reachable on foot, and the child scenes must tile their parent. What fails is returned to the model as notes at most twice, and only what is left after that becomes a gap. The creator's own edits on the plan review page always take precedence and redo what follows from them.

The route is the same loop at every scale, so that shapes are fixed before detail. A parent scene receives its concept, its plan and a rough model; its plan is then split into the shapes of its child scenes, and once those shapes are fixed every child scene is made in parallel, each repeating the loop. The two stages that carry this, \texttt{rough-model} and \texttt{shell-split}, are built but have not yet run.

\subsection{Layers from row kinds}\label{layers-from-row-kinds}

A scene's layers are read off its design alone. Each inventory row's kind, found from the noun the reading gave it by a single table of rules, says what its objects do: emit, move, lie loose, glow, grow, act or sound. A scene receives an effects, motion, physics, crops or people layer only when it has rows that do so, and always receives light and sound. Each layer is made from its rows by the rule for their kind, so a new environment needs no hand-written list of what to build. The scene record holds what no made object covers (shells, stairs, ground and water, backdrop, lamps and neighbouring buildings), and it is gated on the shell lying within 5 cm of the plan and on every doorway opening onto a floor.

Models are made by the prop route first. Nano Banana Pro draws a close-up of each object from its crop of the concept, redrawn at most three times when the shape check rejects it; Pixal3D \citep{li2026pixal3d} builds the model from that close-up; and a likeness gate compares the model's outline from eight sides with the close-up. Only an object whose checks fail is built instead by a code builder that a coding agent writes, and a thin straight object (a pipe, post or beam) whose every close-up has failed is swept along its length in code.

Surfaces come from the creator's library. Everything a scene shows is drawn from score\_lib or recorded as a written exception, and a model contributes its shape only. Each of its parts takes one surface from a library of rule-based surfaces, written as ProcFunc functions \citep{raistrick2026procfunc} on Infinigen's shaders \citep{raistrick2023infinigen} and coloured from the scene's palette. The finish route labels each face from free renders of the model, matches each finish to a library surface, bakes it, and checks the bake against the close-up, requiring each finish to lie within ΔE 32 of its colour there. A free library check then traces every material, model, light, sound and character of an accepted assembly back to the library or to an exception.

\subsection{Settling and physics}\label{settling-and-physics}

Loose objects are dropped onto the made ground, and the settle writes where each comes to rest as a layer of its own; an object that would turn by more than 15 degrees or drift by more than 30 cm is recorded as a resting fault and keeps its laid pose. The physics trial then steps the composed scene on Newton, using MuJoCo's solver on Newton's contacts under the gravity of the body the scene stands on, and judges it by four gates: \texttt{reacted}, \texttt{at\ rest}, \texttt{energy}, which requires that undriven energy does not rise, and \texttt{hinge}, which requires a hinged leaf to swing to its limit.

The trial moved from Blender's Bullet \citep{bullet} to Newton on 27 October 2026, and the move was checked on a toy scene at Moon gravity with two crate drops, a driven door with a wheel fixed to it and a pushed stool. Both engines gave the same answer on every gate, and Newton landed the crates on their targets where Bullet left them on its 8 mm contact margins. Newton's deterministic mode was bitwise equal over six runs on that toy, but it cannot hold scenes with many loose bodies, because its buffers then pass the 32-bit limit; the trial and the settle therefore run without it, and two runs of a 25-body toy agree within 5e-7 m on every body except one that falls off the edge of the floor. The settle is moving to Newton in the same way, so that one engine steps both what comes to rest and what is tried, and one effect is to be stepped on Newton as a first case for the effects layer.

\subsection{Light, motion, effects and sound}\label{light-motion-effects-and-sound}

The sky is computed from the EC's physics rather than drawn, with each body placed at its true direction, size and phase as seen from the scene's site. The lamps are the inventory's glowing rows: one render per scene gives each lamp its own light group, and a free sweep mixes their strengths until the scene lies within 1.5 times its concept's brightness. Every object of a moving row moves by the rule for its kind, at a speed cited from real machines, and is gated on nothing passing within 3 cm of anything else. Emitters are measured against real footage, and each sound source receives either a library file or one MOSS-SoundEffect \citep{moss2026} take, gated on a CLAP \citep{wu2022clap} score of at least 0.30.

\subsection{People and sim bodies}\label{people-and-sim-bodies}

People are made like any other model, from a close-up per person row that Pixal3D builds and the finish route paints. SOMA's body \citep{saito2026soma}, read from the close-up by SAM 3D Body \citep{yang2026sam3dbody}, is fitted beneath the model as its rig, and Kimodo's clips \citep{rempe2026kimodo} move it; a shell gate passes a look when 99 percent of its soft points lie within 5 cm of the body. Each person stands on a person spot of its scene's plan.

Each character also receives a sim body, its physical twin on Newton, so that the people in an environment are simulated and not only animated. The minimal body is built directly from SOMA's 23 segments, with a ball joint at each SOMA joint frame, and is driven by the character's SOMA segment rotations without scaling or per-person setup. Its mass is distributed by segment volume, and it stands on a SCORE floor, so its ground forces come from contact; muscles enter as a minimal first step. Built in this way, the body follows SOMA's rotations to within 6.9e-9 m, and on Newton it matches plain MuJoCo to within 3.4e-8 after 200 steps.

\begin{center}\fbox{\begin{minipage}{0.94\linewidth}\setlength{\parskip}{4pt}\small

\textbf{Gap: Newton for the settle, one effect and the sim body.} The settle and the physics trial on one engine, with the resting and physics gates per scene against the Bullet settle; the one effect stepped on Newton; and a sim body walking on a SCORE floor, its joint angles and ground forces against OpenCap's LabValidation motion capture and force plates \citep{uhlrich2023opencap}. It waits on these three stages being built and run on a cloud card.

\end{minipage}}\end{center}

\subsection{Completion}\label{completion}

A scene is complete when its completion gate passes on its final assembly, which requires every row to be made or dropped with a reason, nothing to be left unmade, no resting fault, and every layer and every shown scene to be present. A fault the creator decides to ship with becomes an accepted gap, recorded with who decided, when and why; the gate lists every gap by name and is never loosened. A hook stops an agent from claiming that a scene is done while its gate fails, because agents judging their own scenes agree with a measured geometry score only 45.8 percent of the time \citep{li2026lego}. Once a scene is complete, its films are made from scripts the creator has locked.
